\title{Geometry Before Multipliers:\\Operation, Motion, Design, and Physically Meaningful Control Coordinates for Synchronous Machines}
\author{}
\date{25 September 2026}
\hypersetup{pdftitle={Geometry Before Multipliers: Operation, Motion, Design, and Physically Meaningful Control Coordinates for Synchronous Machines}}

\newcommand{\PaperAbstract}{%
This whitepaper develops a common geometric solution language for minimum-loss
current allocation and maximum available torque in linear electrically excited
and permanent-magnet synchronous machines. Copper loss, torque, and squared
stator voltage are first classified as quadratic forms. Loss whitening turns
the thermal ellipsoid into a sphere; a physically motivated rotation then
separates torque-producing flux, quadrature current, and a torque-neutral
degree of freedom. This exposes canonical hyperbolic, projective, conformal,
and performance-space descriptions of the same structure. The normalized
performance set is convex in the three-current EESM case, so both operating
problems reduce to a one-dimensional boundary search supported by small
symmetric eigenproblems. The PSM appears as a fixed-excitation slice; its
resistive voltage ellipse, center trajectory, rotation, and characteristic
current cases are derived without prematurely discarding stator resistance.
Resultants, KKT equations, and known PMSM quartics are retained as validation
and active-set tools rather than used as the opening argument. Symbolic and
random numerical tests verify all transformations. The resulting controller
architecture returns feasible minimum-loss currents when possible and the
maximum admissible torque otherwise, while preserving the copper-loss ceiling
as an explicit physical feasibility limit. The same value functions are then
read in reverse: a required torque-speed envelope defines an admissible region
in electromagnetic machine-parameter space. For the loss-limited linear EESM,
two loss-normalized machine coordinates reduce constant torque-per-loss
capability to circles, while their polar angle separates reluctance and field
excitation mechanisms. The singular limit of vanishing excitation resistance
is exposed as a loss cylinder and as a model diagnostic, not confused with the
fixed-excitation PSM slice. Closed design inequalities are derived for the
nonsalient PSM and its characteristic-current window; voltage-limited and
multi-point requirements remain explicit value-function level sets. Spectral
descriptors, sensitivities, robust parameter paths, and a non-unique preimage
from electromagnetic parameters to physical geometry connect the analysis to
machine-control co-design and inexpensive pre-filtering before finite-element
optimization. A third reading asks how the machine should move between two
operating points. Starting from a reciprocal flux-linkage matrix, the full
stator and rotor-excitation dynamics are written as a drift plus a
voltage-controlled velocity. The affine image of the converter set is the
local speedometer of current space. Dynamic whitening separates voltage
authority from electrical drift and identifies fast, medium, and slow current
directions. Minimum time, transient loss energy, actuation effort, and thermal
stress induce different trajectory geometries; their well-posedness conditions
are stated explicitly. Ellipsoidal control authority with sufficiently weak
drift yields an exact Zermelo navigation problem and a Randers metric, whereas
hexagonal stator voltage and an independent excitation converter require a
more general, often polyhedral, Finsler gauge. Curvature, optical refraction,
piecewise-affine magnetic maps, controllability Gramians, reachability
wavefronts, Hamilton--Jacobi equations, and flatness parametrizations are
compared by the physical question each one simplifies. Torque-normal and
torque-tangential motion expose a practical two-stage EESM strategy: acquire
torque with fast stator currents and then redistribute current on the torque
surface while the slower excitation state approaches its loss-optimal value.
A reproducible linear toy experiment quantifies straight, stator-first,
excitation-first, geometrically constructed, time-oriented, and energy-oriented
paths. The resulting transition times, dynamic metric spectra, Gramian
condition numbers, and representative point-to-point distances become new
machine-control co-design descriptors. A fourth reading then asks in which
coordinates the machine should be represented and controlled. Spatial
electromagnetic frames, abstract task coordinates, and virtual input
coordinates are separated rigorously. Rotor-oriented $dq$, stator-flux
$ft$, and active-flux descriptions are derived in both directions, including
the connection terms of moving frames. Torque is shown to be a valid local
coordinate at regular points; hyperbolic torque-allocation coordinates,
performance coordinates $(M,P_{\mathrm{Cu}},U^2)$, natural-gradient frames,
and feedback-linearized task inputs are compared by invertibility,
conditioning, observer burden, and real-time cost. Four numerical machine
cases and symbolic round-trip checks expose singular sets and show why no
single frame is universally best. The resulting conclusion is deliberately
conditional: $dq$ remains an excellent spatial interface, while optimization,
trajectory planning, and EESM allocation often become clearer in different,
task-specific coordinates.}

% Disable either component if the paper does not need it.
\newif\ifshowglossary
\showglossarytrue
\newif\ifshowbibliography
\showbibliographytrue
